% Source: source/普通高考/2016/2016天津文.pdf, page 1, question 11.
% pdfimages -list found no embedded image objects; the program flowchart is vector artwork.
% Grid evidence: tmp/redraw_flowchart_2016_tianjin_liberal/recheck_source/page1_grid100.jpg;
% comparison source is cropped from the no-grid page render.
% Elements: rounded start/end nodes; S=4, n=1, S=2S, S=S-6, n=n+1 process boxes;
% decision diamonds S>=6? and n>3?; output S parallelogram; directed connectors;
% branch labels 是/否 and the right-side loop-back path.
% Relations: start -> S=4 -> n=1 -> S>=6?; yes -> S=S-6 and no -> S=2S;
% both branch results -> n=n+1 -> n>3?; yes -> output S -> end; no -> S>=6?.
% Node-edge list: start->setS; setS->setn; setn->testS; testS(是)->sub;
% testS(否)->double; sub->inc; double->inc; inc->testn; testn(是)->output;
% testn(否)->testS; output->finish. No other connectors are present.
% solid segments: all node borders and connector shafts; dashed segments: none.
% labels: formulas centered in nodes; 是/否 labels beside corresponding exits;
% loop-back arrow enters the top of the S>=6? diamond and avoids the output path.

\begin{tikzpicture}[
    >=Stealth,
    x=1cm,
    y=0.88cm,
    line width=0.45pt,
    line cap=round,
    line join=round,
    font=\normalsize,
    every node/.style={anchor=center, align=center,
        execute at begin node={\setlength{\parindent}{0pt}}},
    flowbox/.style={draw, anchor=center, minimum height=0.68cm, inner xsep=4pt, inner ysep=2pt, align=center,
        execute at begin node={\setlength{\parindent}{0pt}}},
    terminal/.style={flowbox, rounded corners=0.14cm, minimum width=1.25cm},
    io/.style={flowbox, trapezium, minimum height=0.76cm,
        minimum width=1.45cm, trapezium left angle=70, trapezium right angle=110},
    process/.style={flowbox, minimum width=2.65cm},
    decision/.style={flowbox, diamond, aspect=2.7, minimum width=3.35cm,
        minimum height=1.04cm},
]
    % White margins preserve both the left branch and right loop-back.
    \path[use as bounding box] (-3.25,-0.48) rectangle (3.55,11.42);

    \node[terminal] (start) at (0,11.0) {开始};
    \node[process] (setS) at (0,9.72) {$S=4$};
    \node[process] (setn) at (0,8.48) {$n=1$};
    \node[decision, minimum width=3.65cm, minimum height=1.10cm] (testS)
        at (0,7.02) {$S\geq6?$};
    \node[process, minimum width=2.45cm] (double) at (-1.55,5.42) {$S=2S$};
    \node[process, minimum width=2.65cm] (sub) at (0,5.42) {$S=S-6$};
    \node[process, minimum width=2.85cm] (inc) at (0,4.02) {$n=n+1$};
    \node[decision, minimum width=3.35cm, minimum height=1.02cm] (testn)
        at (0,2.58) {$n>3?$};
    \node[io, minimum width=1.45cm] (output) at (0,1.18) {输出 $S$};
    \node[terminal] (finish) at (0,-0.08) {结束};

    % Main chain and the S>=6? yes branch.
    \draw[->] (start.south) -- (setS.north);
    \draw[->] (setS.south) -- (setn.north);
    \draw[->] (setn.south) -- (testS.north);
    \draw[->] (testS.south) -- node[right] {是} (sub.north);
    \draw[->] (sub.south) -- (inc.north);

    % S>=6? no branch is separate from the yes path and merges at n=n+1.
    \coordinate (noTurn) at (-1.55,7.02);
    \coordinate (noDrop) at (-1.55,5.95);
    \draw (testS.west) -- node[above left] {否} (noTurn) -- (noDrop);
    \draw[->] (noDrop) -- (double.north);
    \draw[->] (double.south) |- (inc.west);

    % Increment and second decision.
    \draw[->] (inc.south) -- (testn.north);
    \draw[->] (testn.south) -- node[right] {是} (output.north);
    \draw[->] (output.south) -- (finish.north);

    % n>3? no branch loops back to the first decision without touching output.
    \coordinate (loopRight) at (2.90,2.58);
    \coordinate (loopTop) at (2.90,7.90);
    \coordinate (loopJoin) at (0,7.90);
    \draw (testn.east) -- node[above] {否} (loopRight) -- (loopTop);
    \draw[->] (loopTop) -- (loopJoin) -- (testS.north);
\end{tikzpicture}
